\section{树与生成树}
\label{sec:04-Discrete-Mathematics/03-Graph-Theory/03}

上一节末尾算出一个数：\(n\) 个顶点要连成一片，至少需要 \(n-1\) 条边。把它当成预算来看——你要把公司的 \(n\) 个办公点接进同一张内网，每条专线按年付费，任意两点之间通得了就行，不要求直连。预算压到底，就是恰好铺 \(n-1\) 条。

这时候网络会变成什么样子？没有任何一条线是多余的，因为少一条就有办公点掉出去；也没有任何两点之间存在两条不同的走法，因为多一条走法意味着某条线可以省掉。这两句话描述的是同一种结构，它有个名字：树。

\subsection{加一条成环，减一条断开}

树是连通且无环的无向图。这两个词各自堵住一个方向：连通保证任意两点之间至少有一条路，无环保证至多有一条。若某两点之间有两条不同的路径，把它们拼起来就得到一个环，与无环矛盾。所以树里的路径不仅存在，而且唯一。

唯一性一旦确立，临界性就跟着出来了。

\begin{center}
\begin{tikzpicture}[x=1cm,y=1cm,font=\small,>=stealth]
  % 左panel：加边成环
  \draw[orange!75,line width=2.6pt,opacity=0.55] (-0.8,0) -- (-1.6,1) -- (0,2) -- (1.6,1) -- (0.9,0);
  \draw[black!60] (0,2) -- (-1.6,1);
  \draw[black!60] (0,2) -- (1.6,1);
  \draw[black!60] (-1.6,1) -- (-2.4,0);
  \draw[black!60] (-1.6,1) -- (-0.8,0);
  \draw[black!60] (1.6,1) -- (0.9,0);
  \draw[black!60] (1.6,1) -- (2.4,0);
  \draw[red!75,dashed,thick] (-0.8,0) -- (0.9,0);
  \foreach \x/\y in {0/2,-1.6/1,1.6/1,-2.4/0,-0.8/0,0.9/0,2.4/0}
    {\fill (\x,\y) circle (2.3pt);}
  \node[above] at (0,2.14) {\(r\)};
  \node[left] at (-1.74,1.08) {\(a\)};
  \node[right] at (1.74,1.08) {\(b\)};
  \node[below] at (-0.8,-0.16) {\(d\)};
  \node[below] at (0.9,-0.16) {\(e\)};
  \node[below,red!75,font=\footnotesize] at (0.05,-0.55) {新边 \(de\)};
  \node[font=\footnotesize,black!70] at (0,2.9) {加一条边};
  % 右panel：删边断开
  \begin{scope}[shift={(7.4,0)}]
    \draw[black!40,dashed,fill=blue!5] (-1.2,0.75) ellipse (2.0 and 1.8);
    \draw[black!40,dashed,fill=green!6] (1.6,0.45) ellipse (1.25 and 1.25);
    \draw[black!60] (0,2) -- (-1.6,1);
    \draw[black!25] (0,2) -- (1.6,1);
    \draw[black!60] (-1.6,1) -- (-2.4,0);
    \draw[black!60] (-1.6,1) -- (-0.8,0);
    \draw[black!60] (1.6,1) -- (0.9,0);
    \draw[black!60] (1.6,1) -- (2.4,0);
    \draw[red!75,thick] (0.55,1.72) -- (1.05,1.28);
    \draw[red!75,thick] (0.55,1.28) -- (1.05,1.72);
    \foreach \x/\y in {0/2,-1.6/1,1.6/1,-2.4/0,-0.8/0,0.9/0,2.4/0}
      {\fill (\x,\y) circle (2.3pt);}
    \node[above] at (0,2.14) {\(r\)};
    \node[left] at (-1.74,1.08) {\(a\)};
    \node[right] at (1.74,1.08) {\(b\)};
    \node[font=\footnotesize,black!70] at (0,2.9) {删一条边};
  \end{scope}
\end{tikzpicture}

\emph{同一棵树的两种扰动。左边：新边 \(de\) 的两端本来已有唯一路径 \(d\text{--}a\text{--}r\text{--}b\text{--}e\)（橙色），新边把它闭合成环，而且只闭合出这一个环。右边：删掉 \(rb\)，右侧三个顶点与其余部分再无任何走法。}
\end{center}

图上这两件事各自可以反过来当定义用。``删任何一条边都断开''说的是每条边都不可替代，也就是无环；``加任何一条边都出环''说的是任意两点本来就有路，也就是连通。绕了一圈，回到原地。

于是边数 \(n-1\) 就不是巧合，而是同一件事的第三种说法。

\begin{theorem}[树的三条件]
设 \(G\) 是 \(n\) 个顶点的简单图。连通、无环、恰有 \(n-1\) 条边——这三个条件中任意两条成立，第三条自动成立。
\end{theorem}

不必六条命题两两互推，只要跑通一个方向，其余方向靠反证补齐。骨架是这样的：先证任何至少两个顶点的树都有叶子（度为 1 的顶点）。取一条最长的路径，它的端点若还有别的邻居，这个邻居要么在路径外——那路径还能延长，与最长矛盾；要么在路径上——那就凑出一个环，与无环矛盾。所以端点只有一个邻居。

叶子的用处在于它可以被剥掉：删去一片叶子和它那条边，连通性不受影响（叶子从来不是别人路上的中转站），无环性也不受影响，剩下的还是树。于是 \(n\) 个顶点的树按顶点数归纳，每步剥一片叶子，从一个顶点零条边起算，边数一路加到 \(n-1\)。归纳法为什么能这样用，见\hyperref[sec:01-Mathematical-Language/04-Induction-and-Recursion/01]{``数学归纳法为何成立''}。

剩下两个方向的论证都很短。连通且有 \(n-1\) 条边的图若含环，删掉环上一条边不影响连通，边数降到 \(n-2\)，与上一节``连通至少 \(n-1\) 条边''冲突。无环且有 \(n-1\) 条边的图若不连通，各个分量都是树，边数之和是 \(n\) 减去分量数，小于 \(n-1\)。两处都只用了同一个计数，没有新东西。

\subsection{从任意网络里抠出骨架}

真实的网络当然不会恰好是树。它有环，有冗余，有一堆可省的边。生成树就是在原图里挑出一个子图，顶点一个不少，边只留 \(n-1\) 条，仍然连通。

构造它有两个方向，都只用到上面那张图。往下削：只要还有环，就在环上删一条边——环上剩下的部分照样把两个端点连起来，连通性纹丝不动——删到无环为止。往上长：从一个顶点起步，每次挑一条通向尚未纳入的顶点的边加进来，加到所有顶点都进来为止，这样加的边永远不会闭合成环。两条路都停在同一类结构上，因为 \(n-1\) 条边的连通图只能是树。

原图若本来就不连通，对每个分量各取一棵，得到生成森林。

\subsection{边有价格时，贪心为什么不出错}

给每条边标上年费，不同的生成树总价不同，最便宜的那棵叫最小生成树。这是一个组合优化问题，可行解的数量随 \(n\) 爆炸，然而贪心算法在这里居然是对的——原因还是那张图。

把顶点集任意切成两半，跨越这一刀的边构成一个割。断言是：跨割的最轻边一定出现在某棵最小生成树里。

\begin{center}
\begin{tikzpicture}[x=1cm,y=1cm,font=\small,>=stealth]
  \draw[black!45,dashed] (-0.1,-1.9) -- (-0.1,2.0);
  \node[black!60,font=\footnotesize] at (-0.1,2.3) {割};
  % 当前树的边
  \draw[black!70,line width=1.6pt] (-2.0,1.2) -- (-2.6,0);
  \draw[black!70,line width=1.6pt] (-2.6,0) -- (-1.4,-1.0);
  \draw[black!70,line width=1.6pt] (1.6,1.3) -- (2.4,0.1);
  \draw[black!70,line width=1.6pt] (2.4,0.1) -- (1.3,-1.0);
  \draw[black!70,line width=1.6pt] (-2.0,1.2) -- (1.6,1.3);
  \node[above,black!75,font=\footnotesize] at (-0.2,1.3) {\(f\)：权 \(7\)};
  % 轻边
  \draw[blue!75,dashed,line width=1.6pt] (-1.4,-1.0) -- (1.3,-1.0);
  \node[below,blue!75,font=\footnotesize] at (-0.05,-1.12) {\(e\)：权 \(2\)};
  \foreach \x/\y in {-2.0/1.2,-2.6/0,-1.4/-1.0,1.6/1.3,2.4/0.1,1.3/-1.0}
    {\fill (\x,\y) circle (2.4pt);}
  \node[left] at (-2.72,0) {\(u\)};
  \node[right] at (2.52,0.1) {\(v\)};
\end{tikzpicture}

\emph{粗线是当前的一棵生成树，它用 \(f\) 跨越割线。把更轻的 \(e\) 加进去必然成环，而环从割的一侧出发还要回来，跨割次数是偶数，所以环上除 \(e\) 之外至少还有一条跨割边——图里就是 \(f\)。删 \(f\) 补 \(e\)，仍是生成树，总权不增。}
\end{center}

这个交换论证是最小生成树全部算法的地基。Kruskal 把边按权升序扫一遍，当前边若连接两个尚未连通的分量就收下——它正是这两团顶点划出的那个割上的最轻边；两端已经同属一团就跳过，收下只会成环。Prim 从一个顶点起步维护一棵在长大的树，每步取连接树内与树外的最轻边——树内顶点与树外顶点之间那一刀就是它用的割。两个算法的数据结构和复杂度不同，正确性用的是同一条性质。

顺带一提，Kruskal 扫到一半时的中间状态本身就有意义：按边权从小到大合并，每合并一次就是把两个簇并成一个，这正是单链接层次聚类的过程。所谓聚类树状图，画的就是这条合并轨迹。

\subsection{它没有承诺的事}

最小生成树最小化的是边权之和，不是任何一对顶点之间的距离。树上从 \(u\) 到 \(v\) 只有唯一一条路，这条路可能比原图里的最短路绕出很远——环全被删了，抄近道的备选一并消失。要压缩点对之间的距离，该找的是最短路径树；要在若干指定顶点之间省成本、允许引入中转点，那是 Steiner 树，且它是难解的。名字里都有``树''，目标函数完全不同。

另一件事更要紧：树的零冗余是双向的。省下来的每一分钱，换的是任意一条链路断掉就分裂成两半。上一节的割边概念在这里到达极端——树的每一条边都是割边。真要做容错，就得往骨架上加回备援边，而加回的每一条边恰好制造一个环，那个环就是绕行路线。成本与冗余之间的换算，说到底是环的个数 \(\abs{E}-n+1\)。

树把连通性压到了临界。下一节换一类结构约束：顶点被分成两组，边只在组间出现。这种图上要问的不再是怎样连通，而是怎样配对。
